\bta{2015}


\begin{enumerate}
	%\renewcommand{\labelenumi}{\arabic{enumi}.}
	% A(\Alph) a(\alph) I(\Roman) i(\roman) 1(\arabic)
	%设定全局标号series=example	%引用全局变量resume=example
	%[topsep=-0.3em,parsep=-0.3em,itemsep=-0.3em,partopsep=-0.3em]
	%可使用leftmargin调整列表环境左边的空白长度 [leftmargin=0em]
	\item
	选择题 (本题满分 50 分，每小题 5 分)
	\begin{enumerate}
		%\renewcommand{\labelenumi}{\arabic{enumi}.}
		% A(\Alph) a(\alph) I(\Roman) i(\roman) 1(\arabic)
		%设定全局标号series=example	%引用全局变量resume=example
		%[topsep=-0.3em,parsep=-0.3em,itemsep=-0.3em,partopsep=-0.3em]
		%可使用leftmargin调整列表环境左边的空白长度 [leftmargin=0em]
		\item
		对于函数 $f(x)=\frac{\sin x^{2}}{x}$, 结论不正确的是 \xzanswer{B}


\fourchoices
{在 $(0, \infty)$ 内有界}
{在 $(0, \infty)$ 内 $f(x)$ 没有最大值和最小值}
{在 $(0, \infty)$ 内处处可导}
{当 $x \rightarrow \infty, x \rightarrow 0^{+}$时, $f(x)$ 极限存在}

\item
$ \lim\limits _{x \rightarrow 0}\left(\frac{\sin x}{x}\right)^{\frac{1}{x^{2}}}= $ \xzanswer{D}

\fourchoices
{$ 1 $}
{$ 0 $}
{$\infty$}
{$e^{-\frac{1}{6}}$}

\item
微分方程 $y^{\prime}=\frac{1}{y-x}$ 的通解为 \xzanswer{A}


\fourchoices
{$x=y+C e^{-y}-1$}
{$y=x+C e^{-x}-1$}
{$x=\ln (x-y-1)+C$}
{$y=\ln |y-x-1|+C$}

\item
已知 $m, n$ 为正整数,且 $m>n$ 。如果: $S=\int_{0}^{\frac{\pi}{4}} \sin ^{m} x \cos ^{n} x d x, T=\int_{0}^{\frac{\pi}{4}} \sin ^{n} x \cos ^{m} x d x$。 则下面结论正确的一个是 \xzanswer{C}


\fourchoices
{$S>T$}
{$S=T$}
{$S<T$}
{$S, T$ 的关系不确定}

\item
设对任意的 $x \in \mathbb{R}$, 总有 $m \leqslant f(x)<g(x)<h(x) \leqslant M$ 。 且 $g(x)$ 为连续函数, 若 $\lim\limits _{x \rightarrow \infty}[M-f(x)][h(x)-m]=0$ 。 则对于 $\lim\limits _{x \rightarrow \infty}g(x)$，则下面结论正确的一个是 \xzanswer{B}


\fourchoices
{一定存在, 且等于 $\frac{M+m}{2}$}
{一定存在, 且只能等于 $M$ 或 $m$}
{一定不存在}
{一定存在, 且可以取到 $[m, M]$ 上的任意值}


\item
设函数 $f(x)=x^{2} \sin x+\cos x+\frac{\pi}{2} x$, 在其定义域内零点的个数是 \xzanswer{D}


\fourchoices
{$ 2 $}
{$ 3 $}
{$ 4 $}
{多于 $ 4 $}

\item
设 $ \vv{a} ,  \vv{b} $ 为非零实向量, 且 $2  \vv{a} + \vv{b} $ 与 $ \vv{a} - \vv{b} $ 垂直，$  \vv{a}+2 \vv{b}   $与$  \vv{a} +  \vv{b}  $垂直。则 \xzanswer{\text{条件矛盾}}


\fourchoices
{$| \vv{b} |^{2}=7| \vv{a} |^{2}$}
{$| \vv{a} |^{2}=7| \vv{b} |^{2}$}
{$| \vv{b} |^{2}=5| \vv{a} |^{2}$}
{$| \vv{a} |^{2}=5| \vv{b} |^{2}$}

\item
设平面$ D $由 $x+y=\frac{1}{2}, x+y=1$ 及两条坐标轴围成， $I_{1}=\iint_{D} \ln (x+y)^{3} d x d y, I_{2}=\iint_{D}(x+y)^{3} d x d y, I_{3}=\iint_{D} \sin (x+y)^{3} d x d y$ 。则下面结论正确的一个是 \xzanswer{C}


\fourchoices
{$I_{1}<I_{2}<I_{3}$}
{$I_{3}<I_{1}<I_{2}$}
{$I_{1}<I_{3}<I_{2}$}
{$I_{3}<I_{2}<I_{1}$}


\item
幂级数 $\sum\limits_{n=1}^{\infty} a_{n} x^{n}$ 与 $\sum\limits_{n=1}^{\infty} b_{n} x^{n}$ 的收敛半径分别为 $\frac{\sqrt{2}}{3}$ 与 $\frac{1}{3}$, 则幂级数 $\sum\limits_{n=1}^{\infty} \frac{a_{n}^{2}}{b_{n}^{2}} x^{n}$ 的 收敛半径为 \xzanswer{\text{条件不足}}


\fourchoices
{$ 2 $}
{$\frac{\sqrt{2}}{3}$}
{$\frac{1}{3}$}
{$\frac{1}{2}$}


\item
已知曲面 $x^{2}+y^{2}+2 z^{2}=5$, 在其点 $\left(x_{0}, y_{0}, z_{0}\right)$ 处的切平面与平面 $x+2 y+z=0$ 平行, 则有 \xzanswer{B}


\fourchoices
{$x_{0}: y_{0}: z_{0}=4: 2: 1$}
{$x_{0}: y_{0}: z_{0}=2: 4: 1$}
{$x_{0}: y_{0}: z_{0}=1: 4: 2$}
{$x_{0}: y_{0}: z_{0}=1: 2: 4$}





	\end{enumerate}


\banswer{
\begin{enumerate}
	%\renewcommand{\labelenumi}{\arabic{enumi}.}
	% A(\Alph) a(\alph) I(\Roman) i(\roman) 1(\arabic)
	\item
	\textbf{B}（本题问“不正确”的）。$f$ 是初等函数，在 $(0,+\infty)$ 内处处可导，C 对。
	$\lim\limits_{x\to0^{+}}\frac{\sin x^2}{x}=\lim\limits_{x\to0^{+}}\frac{x^2}{x}=0$，
	而 $|f(x)|\leqslant\frac1x\to0\ (x\to+\infty)$，故 D 对；
	两端极限存在又在内部连续，故 $f$ 有界，A 对。

	B 不正确：取 $x_1=\sqrt{\frac{\pi}{2}}$ 得 $f\left(x_1\right)=\frac{1}{x_1}>0$，
	取 $x_2=\sqrt{\frac{3\pi}{2}}$ 得 $f\left(x_2\right)=-\frac{1}{x_2}<0$。
	由两端极限为 $0$，可取 $0<\alpha<x_1<x_2<\beta$，使 $x\in(0,\alpha)\cup(\beta,+\infty)$ 时
	$|f(x)|<\frac12\min\left\{f\left(x_1\right),\left|f\left(x_2\right)\right|\right\}$。
	$f$ 在闭区间 $[\alpha,\beta]$ 上必取到最大值与最小值，而区间外的值被夹在两者之间，
	故它们就是 $f$ 在 $(0,+\infty)$ 上的最大值与最小值，即最值是存在的。

	\item
	\textbf{D}。由 $\frac{\sin x}{x}=1-\frac{x^2}{6}+o\left(x^2\right)$ 及 $\ln(1+u)\sim u$ 得
	\[
	\frac{1}{x^2}\ln\frac{\sin x}{x}
	=\frac{1}{x^2}\left(-\frac{x^2}{6}+o\left(x^2\right)\right)\to-\frac16 ,
	\]
	故原极限为 $e^{-\frac16}$。

	\item
	\textbf{A}。原方程对 $y$ 不是线性的，但把 $x$ 看成 $y$ 的函数就线性了：
	\[
	\frac{dx}{dy}=y-x\quad\Longrightarrow\quad\frac{dx}{dy}+x=y .
	\]
	乘积分因子 $e^{y}$ 得 $\left(xe^{y}\right)'=ye^{y}$，积分得
	$xe^{y}=(y-1)e^{y}+C$，即 $x=y-1+Ce^{-y}$。

	\item
	\textbf{C}。在 $\left(0,\frac{\pi}{4}\right)$ 上 $0<\sin x<\cos x$，故由 $m>n$ 得
	\[
	\frac{\sin^mx\cos^nx}{\sin^nx\cos^mx}=\tan^{m-n}x<1 ,
	\]
	即 $S$ 的被积函数处处小于 $T$ 的被积函数（且二者连续、不恒等），故 $S<T$。

	\item
	\textbf{B}。由 $f<g<h$ 及 $m\leqslant f$、$h\leqslant M$ 知 $M-f>M-g>0$、$h-m>g-m>0$，故
	\[
	0<\left[M-g(x)\right]\left[g(x)-m\right]<\left[M-f(x)\right]\left[h(x)-m\right]\to0 .
	\]
	记 $u(x)=\left[M-g(x)\right]\left[g(x)-m\right]\to0$，它说明 $g(x)$ 只能越来越靠近 $m$ 或 $M$。
	若 $g$ 既无限次接近 $m$ 又无限次接近 $M$，则由 $g$ 连续及介值定理，
	存在 $z_k\to\infty$ 使 $g\left(z_k\right)=\frac{m+M}{2}$，
	于是 $u\left(z_k\right)=\frac{(M-m)^2}{4}\not\to0$，矛盾。
	故 $\lim\limits_{x\to\infty}g(x)$ 存在且只能是 $M$ 或 $m$。

	\item
	\textbf{D}。定义域为 $\mathbb{R}$，零点有无穷多个。事实上在 $\cos x=0$ 的点上：
	取 $x$ 使 $\sin x=1$，则 $f(x)=x^2+\frac{\pi}{2}x>0\ (x>0)$；
	取 $x$ 使 $\sin x=-1$，则 $f(x)=-x^2+\frac{\pi}{2}x=x\left(\frac{\pi}{2}-x\right)<0\ \left(x>\frac{\pi}{2}\right)$。
	这两类点在 $\left(\frac{\pi}{2},+\infty\right)$ 上交替出现无穷多次，
	由零点定理每两个相邻点之间至少有一个零点，故零点多于 $4$ 个（实为无穷多个）。

	\item
	\textbf{题设矛盾。} 记 $A=|\vv a|^2>0$，$B=|\vv b|^2>0$，$c=\vv a\cdot\vv b$。
    两个垂直条件给出
    \[
    2A-c-B=0,\qquad A+3c+2B=0.
    \]
    代数消元确实得到 $B=7A$ 和 $c=-5A$，但实向量还必须满足
    Cauchy--Schwarz 不等式 $c^2\le AB$。代入后要求
    $25A^2\le7A^2$，与 $A>0$ 矛盾。
    因而不存在符合题意的非零实向量，不能把选项 A 当作有效选择题答案；
    原题某个符号或系数需要核对，现保留题面。

	\item
	\textbf{C}。$D$ 上 $t=x+y\in\left[\frac12,1\right]$，于是 $t^3\in\left[\frac18,1\right]$。
	因 $t\leqslant1$，$\ln t^3=3\ln t\leqslant0$；
	又 $t^3\in\left[\frac18,1\right]\subset\left(0,\frac{\pi}{2}\right)$ 上有 $0<\sin t^3<t^3$。故
	\[
	\ln t^3\leqslant0<\sin t^3<t^3 ,
	\]
	三个被积函数依次递增（且不恒等），积分域相同，所以 $I_1<I_3<I_2$。

	\item
	\textbf{条件不足。} 正确公式是 $R^{-1}=\limsup_{n\to\infty}|c_n|^{1/n}$，
    题设并未保证对应的普通极限存在，也未保证 $b_n\ne0$。
    即使额外要求所有 $b_n\ne0$，结论仍不唯一。例如令
    \[
    a_n=(3/\sqrt2)^n.
    \]
    若 $b_n=3^n$，系数商级数的半径为 $2$。但若
    \[
    b_n=\begin{cases}3^n,&n\text{ 为偶数},\\1,&n\text{ 为奇数},\end{cases}
    \]
    原两级数的半径仍分别为 $\sqrt2/3$ 和 $1/3$，而
    \[
    \limsup_{n\to\infty}|a_n^2/b_n^2|^{1/n}=9/2,
    \]
    所求半径变为 $2/9$。因此不能仅凭题设选择 A。
    若确认原题另有两列根值的极限存在且 $b_n\ne0$ 的条件，则可得 $R=2$。

	\item
	\textbf{B}。曲面 $F=x^2+y^2+2z^2-5$ 在 $\left(x_0,y_0,z_0\right)$ 处的法向为
	$\nabla F=\left(2x_0,2y_0,4z_0\right)$，它应与平面法向 $(1,2,1)$ 平行：
	\[
	\frac{2x_0}{1}=\frac{2y_0}{2}=\frac{4z_0}{1}=\lambda
	\ \Longrightarrow\ x_0=\frac{\lambda}{2},\ y_0=\lambda,\ z_0=\frac{\lambda}{4} ,
	\]
	故 $x_0:y_0:z_0=\frac12:1:\frac14=2:4:1$。
\end{enumerate}

}



\item
(本题满分 10 分)
计算 $\lim\limits _{n \rightarrow \infty} \frac{1}{n}\left(\sqrt{1+\sin \frac{\pi}{n}}+\sqrt{1+\sin \frac{2 \pi}{n}}+\cdots \cdots+\sqrt{1+\sin \frac{n \pi}{n}}\right)$。

\banswer{
	\textbf{化为定积分.} 原式是 $\varphi(x)=\sqrt{1+\sin\pi x}$ 在 $[0,1]$ 上取右端点的
	黎曼和（$\Delta x=\frac1n$，$x_k=\frac kn$），而 $\varphi$ 连续，故
	\[
	\lim_{n\to\infty}\frac1n\sum_{k=1}^{n}\sqrt{1+\sin\frac{k\pi}{n}}
	=\int_0^1\sqrt{1+\sin\pi x}\,dx .
	\]

	\textbf{去根号.} 利用 $1+\sin\theta=\left(\cos\frac{\theta}{2}+\sin\frac{\theta}{2}\right)^2$，
	取 $\theta=\pi x$。当 $x\in[0,1]$ 时 $\frac{\pi x}{2}\in\left[0,\frac{\pi}{2}\right]$，
	正弦余弦均非负，故可直接开方：
	\[
	\sqrt{1+\sin\pi x}=\cos\frac{\pi x}{2}+\sin\frac{\pi x}{2} .
	\]

	\textbf{积分.}
	\[
	\int_0^1\left(\cos\frac{\pi x}{2}+\sin\frac{\pi x}{2}\right)dx
	=\left[\frac{2}{\pi}\sin\frac{\pi x}{2}-\frac{2}{\pi}\cos\frac{\pi x}{2}\right]_0^1
	=\frac{4}{\pi} .
	\]
	即所求极限为 $\frac{4}{\pi}\approx1.2732$（取 $n=10^5$ 的数值和为 $1.27324$，吻合）。
}


\item
(本题满分 10 分)
设 $u=e^{x^{2}} \sin \frac{x}{y}$, 计算 $\frac{\partial^{2} u}{\partial x \partial y}$ 在点 $(\pi, 2)$ 处的值。

\banswer{
	$u$ 在 $y\neq0$ 处有任意阶连续偏导数，混合偏导与次序无关，先对 $y$ 求导较简便：
	\[
	\frac{\partial u}{\partial y}
	=e^{x^2}\cos\frac{x}{y}\cdot\left(-\frac{x}{y^2}\right)
	=-\frac{x}{y^2}e^{x^2}\cos\frac{x}{y} .
	\]
	再对 $x$ 求导（$y$ 视为常数，用乘积法则三项）：
	\begin{align*}
	\frac{\partial^2u}{\partial x\partial y}
	&=-\frac{1}{y^2}\frac{\partial}{\partial x}\left(xe^{x^2}\cos\frac{x}{y}\right)\\
	&=-\frac{e^{x^2}}{y^2}\left[\left(1+2x^2\right)\cos\frac{x}{y}
	-\frac{x}{y}\sin\frac{x}{y}\right] .
	\end{align*}

	\textbf{代入 $(\pi,2)$.} 此时 $\frac{x}{y}=\frac{\pi}{2}$，
	$\cos\frac{\pi}{2}=0$，$\sin\frac{\pi}{2}=1$，故含 $\left(1+2x^2\right)$ 的一项消失：
	\[
	\left.\frac{\partial^2u}{\partial x\partial y}\right|_{(\pi,2)}
	=-\frac{e^{\pi^2}}{4}\left(-\frac{\pi}{2}\right)=\frac{\pi}{8}e^{\pi^2} .
	\]
}



\item
(本题满分 10 分)
设 $D$ 为第一象限内由 $\left\{\begin{array}{l}y=x \\ y=2x \\ x y=1 \\ x y=2\end{array}\right.$所围成的区域, $f$ 为一元可微函数, 且 $f^{\prime}=g$ ，记 $L$ 为 $D$ 的边界。请证明:
\[
\oint_{L} x f\left(\frac{y}{x}\right) \mathrm{d} x=-\int_{1}^{2} \frac{g(u)}{2 u} \mathrm{~d} u
\]


\banswer{
	设 $L$ 取正向（逆时针）。

	\textbf{第一步：格林公式.} 取 $P=xf\left(\frac{y}{x}\right)$、$Q=0$，则
	\[
	\frac{\partial P}{\partial y}=x\cdot f'\left(\frac yx\right)\cdot\frac1x=g\left(\frac yx\right) ,
	\]
	于是
	\[
	\oint_{L}xf\left(\frac yx\right)dx=-\iint_{D}g\left(\frac yx\right)dxdy .
	\]

	\textbf{第二步：换元.} 令 $u=\frac yx$、$v=xy$，则 $D$ 的四条边恰好是
	$u=1$、$u=2$、$v=1$、$v=2$，新区域是正方形 $[1,2]\times[1,2]$。因
	\[
	\frac{\partial(u,v)}{\partial(x,y)}
	=\begin{vmatrix}-\frac{y}{x^2}&\frac1x\\[2pt] y&x\end{vmatrix}
	=-\frac{2y}{x}=-2u ,
	\]
	故 $dxdy=\frac{1}{2u}\,du\,dv$。于是
	\[
	\iint_{D}g\left(\frac yx\right)dxdy
	=\int_1^2\!\!\int_1^2\frac{g(u)}{2u}\,dv\,du=\int_1^2\frac{g(u)}{2u}\,du
	\]
	（对 $v$ 积分的区间长度为 $1$）。代回第一步即得所证。

	\par\vspace{1mm}\noindent
	\textbf{注：}{\kaishu
	雅可比行列式为负，说明 $(u,v)\mapsto(x,y)$ 是反定向的；
	这一点正好与“$L$ 取正向”配套。若把 $L$ 取成顺时针，结论右端要改号。
	取 $f(s)=s^3$（$g=3s^2$）作数值检验，两端同为 $-2.25$。
	}\par
}


\item
(本题满分 10 分)
已知 $y_{1}=x, y_{2}=x^{2}, y_{3}=e^{x}$ 为线性非齐次微分方程: $y^{\prime \prime}+p(x) y^{\prime}+q(x) y=f(x)$ 的 三个特解，求该方程满足初始条件 $y(0)=1, y^{\prime}(0)=0$ 的特解。


\banswer{
	\textbf{构造齐次解.} 两个非齐次特解之差是对应齐次方程的解，故
	\[
	Y_1=y_2-y_1=x^2-x,\qquad Y_2=y_3-y_1=e^{x}-x
	\]
	都满足 $y''+py'+qy=0$。它们线性无关：若
	$\lambda_1\left(x^2-x\right)+\lambda_2\left(e^x-x\right)\equiv0$，
	因 $e^x$ 不是多项式，必有 $\lambda_2=0$，进而 $\lambda_1=0$。

	\textbf{通解.} 二阶线性齐次方程的解空间是二维的，故原方程的通解为
	\[
	y=x+C_1\left(x^2-x\right)+C_2\left(e^{x}-x\right) .
	\]

	\textbf{定常数.} $y(0)=0+0+C_2=1$，故 $C_2=1$；又
	$y'=1+C_1(2x-1)+C_2\left(e^{x}-1\right)$，
	$y'(0)=1-C_1+0=0$，故 $C_1=1$。代回得
	\[
	y=x+\left(x^2-x\right)+\left(e^{x}-x\right)=e^{x}+x^2-x .
	\]
	（验证：$y(0)=1$，$y'=e^x+2x-1$，$y'(0)=0$。）
}


\item
(本题满分 10 分)
设 $f(x)=4 x+\cos \pi x+\frac{1}{1+x^{2}}-x^{2} e^{x}+x e^{x} \int_{x}^{1} f(t) d t$ 。 求 $\int_{0}^{1}(1-x) e^{x} f(x) d x$。

\banswer{
	记所求为 $A=\int_0^1(1-x)e^xf(x)dx$，并记
	\[
	h(x)=4x+\cos\pi x+\frac{1}{1+x^2}-x^2e^{x},\qquad
	F(x)=\int_x^1f(t)dt ,
	\]
	则 $F'(x)=-f(x)$、$F(1)=0$，题设即 $f=h+xe^xF$。

	\textbf{两端积分.} 在 $[0,1]$ 上积分得
	\[
	F(0)=\int_0^1f(x)dx=\int_0^1h(x)dx+\int_0^1xe^{x}F(x)dx .
	\]
	对末项分部积分（注意 $\left[(x-1)e^x\right]'=xe^x$）：
	\begin{align*}
	\int_0^1xe^{x}F\,dx
	&=\Big[(x-1)e^{x}F(x)\Big]_0^1-\int_0^1(x-1)e^{x}F'(x)dx\\
	&=F(0)+\int_0^1(x-1)e^{x}f(x)dx=F(0)-A .
	\end{align*}
	代回上式，$F(0)$ 恰好消去，得到
	\[
	A=\int_0^1h(x)dx .
	\]

	\textbf{算出积分.} 逐项计算：$\int_0^14x\,dx=2$，$\int_0^1\cos\pi x\,dx=0$，
	$\int_0^1\frac{dx}{1+x^2}=\frac{\pi}{4}$，
	$\int_0^1x^2e^{x}dx=\left[e^{x}\left(x^2-2x+2\right)\right]_0^1=e-2$。故
	\[
	A=2+0+\frac{\pi}{4}-(e-2)=4+\frac{\pi}{4}-e\approx2.0671 .
	\]

	\par\vspace{1mm}\noindent
	\textbf{注：}{\kaishu
	本题的巧处在于：所求的 $A$ 与未知的 $F(0)=\int_0^1f$ 通过同一个分部积分联系起来，
	$F(0)$ 自动抵消，因而不必真的解出 $f$。
	}\par
}


\newpage
\item
(本题满分 10 分)
计算曲面积分 $I=\iint_{\Sigma} \frac{x d y d z+y d z d x+z d x d y}{\left(x^{2}+y^{2}+z^{2}\right)^{\frac{3}{2}}}$, 其中 $\Sigma$ 是曲面 $2 x^{2}+2 y^{2}+z^{2}=4$ 的外侧。

\banswer{
	记 $\bm r=(x,y,z)$、$r=|\bm r|$，则被积向量场为 $\bm F=\frac{\bm r}{r^3}$。

	\textbf{第一步：场是无源的.} 在 $\bm r\neq\bm 0$ 处
	\[
	\frac{\partial}{\partial x}\frac{x}{r^3}=\frac{1}{r^3}-\frac{3x^2}{r^5} ,
	\]
	三项相加得 $\operatorname{div}\bm F=\frac{3}{r^3}-\frac{3r^2}{r^5}=0$。
	但原点是奇点且落在椭球 $2x^2+2y^2+z^2=4$ 内部，故不能直接用高斯公式。

	\textbf{第二步：挖去奇点.} 取充分小的 $\varepsilon>0$，作球面
	$\Sigma_\varepsilon:\ x^2+y^2+z^2=\varepsilon^2$（含于 $\Sigma$ 内），取外侧。
	在 $\Sigma$ 与 $\Sigma_\varepsilon$ 夹成的区域 $\Omega$ 上 $\bm F$ 光滑且无源，由高斯公式
	\[
	\iint_{\Sigma}\bm F\cdot d\bm S-\iint_{\Sigma_\varepsilon}\bm F\cdot d\bm S
	=\iiint_{\Omega}\operatorname{div}\bm F\,dV=0 .
	\]

	\textbf{第三步：算小球面.} 在 $\Sigma_\varepsilon$ 上 $r=\varepsilon$，
	外法向单位向量 $\bm n=\frac{\bm r}{\varepsilon}$，故
	\[
	\bm F\cdot\bm n=\frac{\bm r\cdot\bm r}{\varepsilon^3\cdot\varepsilon}
	=\frac{\varepsilon^2}{\varepsilon^4}=\frac{1}{\varepsilon^2} ,
	\]
	于是 $\iint_{\Sigma_\varepsilon}\bm F\cdot d\bm S=\frac{1}{\varepsilon^2}\cdot4\pi\varepsilon^2=4\pi$。
	所以
	\[
	I=4\pi .
	\]

	\par\vspace{1mm}\noindent
	\textbf{注：}{\kaishu
	结果与曲面的具体形状无关，只取决于它是否包围原点：这正是
	$\operatorname{div}\frac{\bm r}{r^3}=4\pi\delta(\bm r)$ 的严格版本。
	直接对椭球面参数化作数值积分，得 $I\approx12.5664=4\pi$。
	}\par
}


\item
(本题满分 10 分)
将函数 $f(x)=x-1$ $(0 \leqslant x \leqslant 2)$ 展开成周期为$  4  $的余弦级数。

\banswer{
	作偶延拓（延拓后为 $|x|-1$，$-2\leqslant x\leqslant2$）再作周期 $4$ 的延拓，此时
	半周期 $l=2$，$b_n=0$，
	\[
	a_n=\frac{2}{l}\int_0^{l}f(x)\cos\frac{n\pi x}{l}dx=\int_0^2(x-1)\cos\frac{n\pi x}{2}dx .
	\]

	\textbf{定 $a_0$.} $a_0=\int_0^2(x-1)dx=0$，故常数项为零。

	\textbf{定 $a_n\ (n\geqslant1)$.} 因 $\int_0^2\cos\frac{n\pi x}{2}dx=0$，只需算 $x$ 那一项：
	\[
	\int_0^2x\cos\frac{n\pi x}{2}dx
	=\left[\frac{2x}{n\pi}\sin\frac{n\pi x}{2}+\frac{4}{n^2\pi^2}\cos\frac{n\pi x}{2}\right]_0^2 ,
	\]
	代入 $\sin n\pi=0$、$\cos n\pi=(-1)^n$ 得 $a_n=\frac{4\left[(-1)^n-1\right]}{n^2\pi^2}$。
	即 $n$ 为偶数时 $a_n=0$，$n=2k+1$ 时 $a_n=-\frac{8}{n^2\pi^2}$。

	\textbf{结论.} 延拓后的函数处处连续，故级数在 $[0,2]$ 上处处收敛于 $f$：
	\begin{align*}
	f(x)&=\left(\frac{2}{\pi}\right)^2\sum_{n=1}^{\infty}\frac{(-1)^n-1}{n^2}\cos\frac{n\pi x}{2}\\
	&=-\frac{8}{\pi^2}\sum_{k=0}^{\infty}\frac{1}{(2k+1)^2}\cos\frac{(2k+1)\pi x}{2} .
	\end{align*}

	\par\vspace{1mm}\noindent
	\textbf{注：}{\kaishu
	取 $x=0$ 得 $-1=-\frac{8}{\pi^2}\sum\limits_{k\geqslant0}\frac{1}{(2k+1)^2}$，
	即 $\sum\limits_{k\geqslant0}\frac{1}{(2k+1)^2}=\frac{\pi^2}{8}$，可作自检；
	在 $x=0.3,\ 0.7,\ 1.5$ 处对级数求和，数值分别为
	$-0.7000,\ -0.3000,\ 0.5000$，与 $x-1$ 一致。
	}\par
}


\item
(本题满分 10 分)
若 $g(x)$ 为单调增加的可微函数, 且当 $x \geqslant a$ 时, $\left|f^{\prime }(x)\right| \leqslant g^{\prime }(x)$ 。 证明：当 $x \geqslant a$ 时, $|f(x)-f(a)| \leqslant g(x)-g(a)$。


\banswer{
	把要证的不等式拆成两个单边不等式，各造一个辅助函数：
	\begin{align*}
	\varphi(x)&=\left[g(x)-g(a)\right]-\left[f(x)-f(a)\right],\\
	\psi(x)&=\left[g(x)-g(a)\right]+\left[f(x)-f(a)\right].
	\end{align*}
	显然 $\varphi(a)=\psi(a)=0$。由 $\left|f'(x)\right|\leqslant g'(x)$ 知
	$-g'(x)\leqslant f'(x)\leqslant g'(x)$，故当 $x\geqslant a$ 时
	\[
	\varphi'(x)=g'(x)-f'(x)\geqslant0,\qquad
	\psi'(x)=g'(x)+f'(x)\geqslant0 .
	\]
	于是 $\varphi$、$\psi$ 在 $[a,+\infty)$ 上单调不减，从而
	$\varphi(x)\geqslant\varphi(a)=0$、$\psi(x)\geqslant\psi(a)=0$，即
	\[
	-\left[g(x)-g(a)\right]\leqslant f(x)-f(a)\leqslant g(x)-g(a) ,
	\]
	也就是 $|f(x)-f(a)|\leqslant g(x)-g(a)$。

	\par\vspace{1mm}\noindent
	\textbf{注：}{\kaishu
	这里只用了可导性与“导数非负则函数不减”，
	既不要求 $f'$ 连续（故不必写成积分），
	也避开了柯西中值定理在 $g'=0$ 处需单独讨论的麻烦。
	}\par
}


\item
(本题满分 10 分)
函数 $f(x)$ 在 $[0,2]$ 上二阶可导, 且 $f^{\prime}(0)=f^{\prime}(2)=0$ 。 证明：在区间 $(0,2)$ 内至少存 在一点 $\xi$, 满足 $\left|f^{\prime \prime}(\xi)\right| \geqslant|f(2)-f(0)|$。

\banswer{
	在 $x=0$ 与 $x=2$ 处分别展成带拉格朗日余项的一阶泰勒公式，并在中点 $x=1$ 处取值。
	注意 $f'(0)=f'(2)=0$，$(1-0)^2=(1-2)^2=1$：
	\begin{align*}
	f(1)&=f(0)+\frac{f''\left(\xi_1\right)}{2},&&\xi_1\in(0,1),\\
	f(1)&=f(2)+\frac{f''\left(\xi_2\right)}{2},&&\xi_2\in(1,2).
	\end{align*}
	两式相减得
	\[
	f(2)-f(0)=\frac{f''\left(\xi_1\right)-f''\left(\xi_2\right)}{2} .
	\]
	取 $\xi$ 为 $\xi_1,\xi_2$ 中使 $\left|f''\right|$ 较大者，则
	\[
	|f(2)-f(0)|\leqslant\frac{\left|f''\left(\xi_1\right)\right|
	+\left|f''\left(\xi_2\right)\right|}{2}\leqslant\left|f''(\xi)\right| ,
	\]
	且 $\xi\in(0,2)$，即为所求。

	\par\vspace{1mm}\noindent
	\textbf{注：}{\kaishu
	这是一般结论 $\left|f''(\xi)\right|\geqslant\frac{4}{(b-a)^2}|f(b)-f(a)|$
	在 $[a,b]=[0,2]$ 时的特例（此时 $\frac{4}{(b-a)^2}=1$）。
	}\par
}



\item
(本题满分 10 分)
设 $0<a<1$，$ x \geqslant 0$ 且 $y \geqslant 0$ 。 证明:
\begin{enumerate}
	%\renewcommand{\labelenumi}{\arabic{enumi}.}
	% A(\Alph) a(\alph) I(\Roman) i(\roman) 1(\arabic)
	%设定全局标号series=example	%引用全局变量resume=example
	%[topsep=-0.3em,parsep=-0.3em,itemsep=-0.3em,partopsep=-0.3em]
	%可使用leftmargin调整列表环境左边的空白长度 [leftmargin=0em]
	\item
$x^{a} y^{1-a} \leqslant a x+(1-a) y$
\item
设 $x_{1}, x_{2}, x_{3}, \cdots, x_{n}, y_{1}, y_{2}, \cdots, y_{n}$, 利用 (1) 中的结果证明:
\[
x_{1} y_{1}+x_{2} y_{2}+\cdots+x_{n} y_{n} \leqslant\left(x_{1}^{2}+x_{2}^{2}+\ldots+x_{n}^{2}\right)^{\frac{1}{2}}\left(y_{1}^{2}+y_{2}^{2}+\ldots+y_{n}^{2}\right)^{\frac{1}{2}}
\]
\end{enumerate}


\banswer{
	\begin{enumerate}
		%\renewcommand{\labelenumi}{\arabic{enumi}.}
		% A(\Alph) a(\alph) I(\Roman) i(\roman) 1(\arabic)
		%设定全局标号series=example	%引用全局变量resume=example
		%[topsep=-0.3em,parsep=-0.3em,itemsep=-0.3em,partopsep=-0.3em]
		%可使用leftmargin调整列表环境左边的空白长度 [leftmargin=0em]
		\item
		若 $x=0$ 或 $y=0$，则左端为 $0$（注意 $0<a<1$），而右端非负，不等式显然成立。
		以下设 $x>0$、$y>0$。

		\textbf{解法一（化为一元）.} 两端同除以 $y$，令 $t=\frac{x}{y}>0$，需证
		$t^a\leqslant at+1-a$。令 $\varphi(t)=at+1-a-t^a$，则 $\varphi(1)=0$，且
		\[
		\varphi'(t)=a-at^{a-1}=a\left(1-t^{a-1}\right) .
		\]
		因 $a-1<0$：$0<t<1$ 时 $t^{a-1}>1$，$\varphi'<0$；$t>1$ 时 $\varphi'>0$。
		故 $t=1$ 是 $\varphi$ 的最小值点，$\varphi(t)\geqslant\varphi(1)=0$，即得所证。

		\textbf{解法二（凹性）.} $(\ln t)''=-\frac{1}{t^2}<0$，$\ln$ 是凹函数，故
		\[
		\ln\left[ax+(1-a)y\right]\geqslant a\ln x+(1-a)\ln y=\ln\left(x^ay^{1-a}\right) ,
		\]
		两端取指数即可。（一般地，$f''\leqslant0$ 时有 $f(ax+by)\geqslant af(x)+bf(y)$，$a+b=1$。）

		\item
		记 $X=\left(\sum\limits_{i=1}^{n}x_i^2\right)^{1/2}$、
		$Y=\left(\sum\limits_{i=1}^{n}y_i^2\right)^{1/2}$。
		若 $X=0$ 或 $Y=0$，则左端为 $0$，结论显然。设 $X,Y>0$。

		在 (1) 中取 $a=\frac12$，并以 $\frac{x_i^2}{X^2}$、$\frac{y_i^2}{Y^2}$ 代替其中的 $x$、$y$：
		\[
		\frac{\left|x_i\right|}{X}\cdot\frac{\left|y_i\right|}{Y}
		\leqslant\frac12\left(\frac{x_i^2}{X^2}+\frac{y_i^2}{Y^2}\right) .
		\]
		对 $i=1,\dots,n$ 求和，右端两个和式各为 $1$：
		\[
		\frac{1}{XY}\sum_{i=1}^{n}\left|x_iy_i\right|\leqslant\frac12(1+1)=1 .
		\]
		故 $\sum\limits_{i=1}^{n}x_iy_i\leqslant\sum\limits_{i=1}^{n}\left|x_iy_i\right|\leqslant XY$，
		即柯西--施瓦茨不等式。


	\end{enumerate}


}



\end{enumerate}

